\begin{helper}[Replacement-pair selection and switch invariance]
Let $C$ be a \noderef{ScheideReducedConfiguration} for
$(G,k,e_0,\varphi,T)$, let $D$ be the associated
\noderef{ScheideDefectiveSetup}, let $R$ assign
\noderef{ScheideReplacementData} to every defective color, and let $L$ be
a \noderef{ScheideLastExitPath} certificate for $\alpha$.  Suppose that
$L.\mathit{lastInside}\notin V(T)$.  There is a defective color $\delta$
such that
\[
 L.\mathit{lastInside}\in\{(R\delta).u1,(R\delta).u2\}.
\]
If $r=D.\gammaIndex(\delta)$ and
$\varphi_\delta$ is the
\noderef{TwoColorComponentSwap} of $\gamma=C.\gamma(r)$ and $\delta$ on
the component rooted at $D.v0(\delta)$, then
\[
 \noderef{TwoColorIncidenceGraph}(G,\varphi_\delta,D.\alpha_0,\alpha)
 =\noderef{TwoColorIncidenceGraph}(G,\varphi,D.\alpha_0,\alpha).
\]
\end{helper}
\begin{proof}
Put $v=L.\mathit{lastInside}$ and $X_0=V(T)$.  The certificate $L$ says
that $v$ belongs to \noderef{ScheideEnlargedVertexSet}$(R)$, whose
definition is
\[
 X_0\cup\bigcup_{d\in\noderef{DefectiveColors}(G,\varphi,X_0)}
       \{(Rd).u1,(Rd).u2\}.
\]
Since $v\notin X_0$, membership in this indexed union supplies a defective
index $\delta$ with $v\in\{(R\delta).u1,(R\delta).u2\}$.

Let $r=D.\gammaIndex(\delta)$ and $\gamma=C.\gamma(r)$.  By
\noderef{ScheideReplacementData}, the appropriate one of
$(R\delta).f1,(R\delta).f2$ is a $\delta$-colored edge incident with $v$;
its other endpoint lies in
$X_0\setminus\{C.\mathit{blockVertex}(r,0),
C.\mathit{blockVertex}(r,1)\}$.  The second conclusion of
\noderef{ScheideSymmetricDefectiveInternal} supplies a $\gamma$-colored
edge between $(R\delta).u1$ and $(R\delta).u2$, hence also incident with
$v$.  Both edges have both endpoints in the enlarged set, whereas the
edge $L.\mathit{crossing}$ leaves that set.  Thus both are distinct from
$L.\mathit{crossing}$.  The consecutive vertex-edge incidence recorded
by the walk in \noderef{ScheideLastExitPath} makes
$L.\mathit{crossing}$ incident with $v$, and its color is $\alpha$.

The extension field of \noderef{ScheideReducedConfiguration} supplies the
\noderef{ScheideStructuralConfiguration}; expanding its
\noderef{GlobalMaximumTriple} and \noderef{KTriple} fields gives
\noderef{IsProperEdgeColoring}$(G,\varphi)$.  Properness and
\noderef{AdjacentEdges}, applied at $v$ to the three distinct incident
edges just identified, give
\[
 \alpha\ne\delta\quad\hbox{and}\quad\alpha\ne\gamma.       \tag{1}
\]

The setup \noderef{ScheideDefectiveSetup} gives
$D.\alpha_0\ne\gamma$.  Since $\delta$ belongs to
\noderef{DefectiveColors}, expansion through \noderef{DefectiveColor}
provides a $\delta$-colored boundary edge of $X_0$.  The theorem
\noderef{ScheideGlobalMaximumOrder} says that $X_0$ is
\noderef{ColoredVertexSetClosed}, while $D.\alpha_0$ is missing at $y_0$
by \noderef{ScheideDefectiveSetup}.  A color missing on a closed set has
no boundary edge, so
\[
 D.\alpha_0\ne\delta\quad\hbox{and}\quad
 D.\alpha_0\ne\gamma.                                    \tag{2}
\]

The formula defining \noderef{TwoColorComponentSwap} changes an edge
color only by interchanging $\gamma$ and $\delta$.  Equations (1)--(2)
therefore give, for every edge identity $q$,
\[
 \varphi_\delta(q)=D.\alpha_0\iff\varphi(q)=D.\alpha_0,
 \qquad
 \varphi_\delta(q)=\alpha\iff\varphi(q)=\alpha.
\]
After expanding \noderef{TwoColorIncidenceGraph}, its adjacency relation
depends only on incidence and these two color tests.  The two graphs are
therefore equal, as claimed.
\end{proof}
